\chapter{The second circular, from source to an unreleased control}
\label{ch:circular}

\section{The source packet and the question it supports}
\label{sec:circular-packet}
The running example begins with an identifiable public document: SFC circular
23EC46, issued on 24 October 2023, concerning additional returns and other
services or arrangements offered when marketing SFC-authorised funds
\citep{sfcmarketing}. The retained experiment selected its paragraph-10 gift
restriction. The circular as a whole discusses more than that restriction. A
reader needs to know both what was selected and what remains outside the control
before a result from the control can be interpreted.

The experiment preserved the official source response, an extracted text
derivative and anchors connecting the selected rule to its source passages.
On 22 September 2026, a fresh retrieval matched the retained raw bytes. This
check establishes that the retained response matched what the retrieval route
returned at that time. It does not establish the complete current legal position
or resolve the effective versions of incorporated sources. A hash identifies
bytes; it does not by itself establish their authority, completeness or legal
effect.

The selected assessment concerns one proposed incentive component, offered by a
distributor in promoting a fund within the declared profile. It is a question
about a specific restriction. It is not a whole-offer compliance verdict. The
distinction is visible in the result vocabulary. A true result means that the
selected prohibition trigger is established. A false result means that this
trigger is not established as true on the supplied, resolved facts under the
specified semantics. It supplies no general permission to market the offer.

The retained packet explicitly leaves the applicable Code paragraph and cited
FAQ unresolved. This matters twice. Those sources may settle the meaning or
scope of a term; they may also establish that the selected profile is narrower
than the underlying requirement. Their absence is therefore a release-blocking
dependency, not a bibliographic inconvenience. The example is useful precisely
because it continues to expose that absence after its executable tests pass.

The book uses this historical, selected packet to teach implementation and
assurance. New commercial offers, cashback definitions and mixed-component
scenarios introduced below are synthetic investigations. They are not reported
SFC rulings. The unresolved references are not silently filled using a language
model's recollection of a Code provision. Doing so would convert a visible gap
into an untraceable premise.

\section{Reading the surrounding circular before isolating the rule}
\label{sec:circular-context}
A translation that begins by copying the most convenient sentence can miss why
that sentence is there. The circular's opening establishes its distribution
context. The next passages describe observed practices, including additional
features, standing investment arrangements and marketing. Those observations
help identify the activity being discussed, but they are not automatically
individual executable prohibitions. A source inventory should retain their
function without manufacturing a rule for every paragraph.

The discussion then turns to additional returns and their possible regulatory
consequences. Some questions concern the structure of an arrangement and public
promotion. Others concern impressions created about the underlying fund. The
gift discussion appears within that broader context. Consequently, a finding
that one separate classification is absent does not clear the gift issue. The
gift control is not conditional on first finding that every other regulatory
concern applies.

Later passages address redemption arrangements, administrative cutoffs,
advertising and the description of services. Some concerns require event or
workflow evidence; others require qualitative assessment of communication.
Those mechanisms differ from the paragraph-10 Boolean control. An adequate
implementation of the whole circular would need to preserve those differences.
It cannot simply accumulate a larger expression around the gift predicate and
call the circular complete.

The existing disposition inventory has 18 numbered paragraphs and four
footnotes. Its assignments are agent-authored and unreviewed. They include
context, separate controls, qualitative review and unresolved dependencies.
Paragraph 10 is the only numbered passage with an executable control in this
exercise. A future independent reviewer should challenge both the assignments
and the completeness of the inventory. Merely approving the final formula would
leave the earlier decision about what to implement unexamined.

This provides an immediate requirement for the ensemble. One path must inspect
the source independently of the candidates and ask what obligations, exceptions
and dependencies need a disposition. Another path may develop the selected
rule. Comparing those paths creates a place for an omitted requirement to
appear. If both start from the same already-filtered list of paragraphs, the
apparent redundancy begins too late.

\section{Naming the six inputs precisely}
\label{sec:circular-inputs}
The executable example has two scope inputs and four body inputs. Their names
are convenient for software, but their definitions do the substantive work.
The first scope input records whether the assessed actor is a distributor in
the declared profile. The second records whether the activity concerns promotion
of an SFC-authorised fund within that profile. Both are reviewed classifications.
The prototype does not infer them from the actor's name or the presence of the
word ``fund'' in a document.

The first body input records whether the proposed component is classified as a
gift. The second and third record a promotional link to a specific product and
to a particular product type. A link to funds from a particular fund house
requires an explicit classification where the relevant example is used; the
engine does not assume that every mention of a fund house creates such a link.
The fourth body input records whether the entire assessed component falls
within the fee-or-charge-discount exception.

The word ``entire'' in that final definition is an engineering expression of
the proposed component-level reading. It prevents an assessed voucher from
inheriting the classification of a separate fee waiver. It does not prove that
the decomposition is legally correct. A reviewer must establish the relevant
unit of assessment from the source and context. The definition and the review
issue must travel together until that question is resolved.

For a compact mathematical account, let $A$ denote established actor scope and
$S$ the fund/activity scope. Let $G$ denote the gift classification, $P$ the
specific-product link, $T$ the product-type link and $D$ the discount exception.
For now suppose these are known Boolean values. The body condition is
\begin{equation}
 B=G\land(P\lor T)\land\neg D.
 \label{eq:gift-body}
\end{equation}
It asks three things together: is the component a gift, does either promotional
route hold, and is the exception absent? The two alternatives occur inside one
parenthesized condition. Replacing that disjunction by a test of $P$ alone
changes the rule for offers linked only to a product category.

Scope is assessed separately. When both scope inputs are true, the body answers
the selected question. When a scope input is known false, the limited profile
returns \texttt{OUT\_OF\_SCOPE}. When scope is unresolved, it returns
\texttt{UNKNOWN}, even if a body condition might otherwise be false. The
distinction prevents an unresolved applicability question from being concealed
inside an apparently definitive in-scope result.

The current runtime also has a conservative conflict policy. It checks the
static dependency closure of the requested rule before evaluating branches.
Conflicting referenced evidence produces \texttt{CONFLICT}, including when a
later scope branch would otherwise make the body irrelevant. This policy is
part of the declared engineering semantics. It should not be attributed to the
wording of the SFC circular. Chapter~\ref{ch:proof} explains why such a policy
must be fixed before comparing implementations.

\section{The named cases are an argument about the rule}
\label{sec:circular-cases}
The 22 named cases are most useful when read as a sequence of challenges to an
interpretation. They do not merely fill rows in a test fixture. Each changes a
fact or evidentiary condition that a plausible implementation might mishandle.
The baseline case is a component classified as a gift, linked to one specific
fund, with no fee-discount exception, and known in-scope actor and activity.

\subsection{The alternative promotional route}
Case g01 establishes the baseline prohibition trigger. Under
equation~\eqref{eq:gift-body}, $G$ and $P$ are true and $D$ is false, so the
body is true. This is a necessary example but a weak test on its own. An
incorrect expression $G\land P$ would also return true, even though it omits
both an exception and an alternative route.

Case g02 changes the promotional facts: no individual fund is named, but the
component is linked to a fund category. Now $P$ is false and $T$ is true.
The disjunction remains true. A candidate that omitted the product-type route
would return false, making this case a concrete witness to the omission. The
source's inclusion of that route supplies the interpretive objection; the
different program outputs make its consequence inspectable.

Case g15 uses the fund-house example through a supplied reviewed product-type
classification. Its purpose is not to write a text classifier for fund-house
mentions. It checks that the chosen formal vocabulary has a place to represent
the classification when review establishes it. This distinction prevents a
source example from being silently generalized into a universal inference from
an entity name.

Case g09 establishes a specific-product link while leaving the product-type
link unknown. The known route is sufficient for the disjunction. Unknown
information about a second route does not erase the first. Case g14 reverses
the situation: the specific-product route is false and the remaining route is
unknown. Here the unresolved fact can change the outcome, so the correct result
under the declared semantics is unknown.

Taken together, these cases explain the meaning of ``either'' in an executable
rule with partial evidence. The system must neither require every route to be
known true nor convert every missing route to false. This is a semantic
distinction, not a preference about how verbose the explanation should be.

\subsection{The exception and the object it qualifies}
Case g03 changes the baseline component to one classified as solely a fee or
charge discount. The exception settles the selected body condition as false.
The result says that this particular gift restriction is not triggered under
the supplied classification. It does not resolve advertising, authorization or
any other requirement associated with the offer.

Case g04 returns to a voucher while another component of the same offer is a
fee waiver. The assessed component remains the voucher, so the waiver does not
set its $D$ value to true. This case challenges an implementation that stores a
single offer-level discount flag and reuses it for every component. If that flag
is the only available data, the adapter has insufficient information to apply
the proposed component-level rule. It should request a finer assessment rather
than derive a convenient value.

Case g07 leaves the discount classification unresolved for a component otherwise
meeting the trigger. If $D$ is false, the prohibition is triggered; if $D$ is
true, the exception applies. Both possibilities remain consistent with the
supplied evidence. Unknown is therefore the informative answer. A system that
turns this unknown into false would remove the exception by default; one that
turns it into true would invent the exception.

Case g16 provides known discount evidence while leaving the gift and promotional
classifications unknown. Within known scope, a true $D$ makes $\neg D$ false,
which settles this conjunction as false. The case demonstrates that every
missing fact need not prevent every narrow conclusion. Its safety depends on
the specific conclusion being reported. It would be unsafe to broaden the
result into a complete clearance of the marketing arrangement.

The contrast between g04 and g16 is particularly useful for reviewer training.
An established exception can settle a condition, but only for the object to
which it applies. A precise value attached to the wrong subject is not better
evidence than an unknown value attached to the right one.

\subsection{Classification is not supplied by the formula}
Case g08 concerns additional returns whose gift classification has not been
settled. The circular's qualified treatment of such returns motivates an
assessment; it does not establish $G$ for every arrangement. The proposed
candidate must retain unknown when the remaining facts would make that
classification decisive. Replacing the missing judgment by an unconditional
``additional return implies gift'' would create a new interpretive rule.

Cases g06 and g22 establish that no gift is offered. In g06 the discount
classification is also missing. The false gift predicate settles the body as
false under the declared Boolean semantics. Case g22 makes the point with a
known product link: a promotional link alone does not establish that a gift
exists. These cases catch a candidate that assumed $G$ from the context instead
of requiring evidence for it.

Case g05 establishes that neither product-link route holds. This also makes the
body false, while leaving other regulatory questions outside the result. It
does not establish that every incentive unrelated to an identified product is
permissible under all applicable requirements. The selected condition is only
one element of a broader assessment.

Case g17 separates regulatory classifications. A separate assessment finds no
structured product, but the component is a non-discount gift with a relevant
product link. The gift trigger remains true. The case challenges a workflow
that exits early after resolving one concern and never evaluates another.
Correct integration therefore requires an inventory of outstanding controls,
not a single reassuring flag named \texttt{regulatory\_check\_passed}.

\subsection{Applicability precedes an in-scope answer}
Cases g10 and g11 establish that, respectively, the actor and the product/activity
are outside the selected profile. The response is out of scope. The underlying
Code requirement may have a broader reach; the example cannot answer that
broader question. Treating out of scope as false would erase the distinction
between an in-scope condition that is not triggered and a question the selected
control does not purport to decide.

Case g12 leaves the fund/activity scope unknown. The system cannot establish
whether the profile applies. Case g13 leaves actor scope unknown while
establishing no gift. A Boolean programmer might be tempted to use the false
gift fact to return false immediately. The declared profile instead reports
unknown scope. This preserves the fact that the system has not established the
context in which its body answer would be meaningful.

This is a design choice that an independent implementation must reproduce.
Another language could define a different scope-combination policy, but it
would then implement a different specification. Differential testing only has
a clear target after the policy has been made explicit. A reviewer must also
decide how an unknown scope result affects the transaction workflow.

\subsection{Contradiction, expiry and later knowledge}
Case g18 supplies conflicting gift evidence. The result is conflict, not
unknown and not the value favored by the latest-arriving record. Unknown means
the evidence does not determine a value. Conflict means admissible records
assert incompatible values. These require different remediation: obtaining
missing evidence in the first case, and resolving or superseding contradictory
evidence in the second.

Case g19 combines an out-of-scope actor with conflicting discount evidence.
The current static-closure conflict policy returns conflict before scope
evaluation. This is deliberately conservative. The example must teach the
policy honestly because a reviewer might otherwise expect out of scope and
mistake the difference for a Java bug. The policy is recorded as engineering
semantics, subject to separate operational review.

Case g20 makes the gift evidence stale at assessment time. The rule does not
continue treating an expired classification as current. Case g21 places the
recording time after the historical knowledge cutoff. A historical assessment
cannot use that later knowledge. Both become unknown when the missing gift
classification remains decisive, but their reason codes and repair actions
differ. Refreshing an expired record does not justify rewriting what was known
at an earlier point in time.

All 22 cases therefore contribute a distinct aspect of the specified behavior.
Some concern the wording of the selected restriction. Others concern evidence,
scope and operational conventions needed to execute the reading. The case
packet should label those sources of expected behavior separately. Otherwise
an implementation convention can acquire the apparent authority of a legal
sentence simply because both are stored in the same JSON file.

\section{What the 729 combinations cover}
\label{sec:circular-enumeration}
Each of the six abstract inputs takes one of three values: true, false or
unknown. The number of assignments is therefore
\begin{equation}
 3^6=729.
 \label{eq:case-domain-size}
\end{equation}
This is a statement about a finite abstraction. It does not count legal
interpretations, marketing documents, evidence histories or real client offers.
The enumeration is exhaustive for that abstraction because it visits every
combination of the six declared values.

The independent truth oracle handles scope first. For the four body facts it
replaces each unknown by both possible Boolean values and evaluates the formula
under every completion. If every completion returns true, the oracle returns
true; if every completion returns false, it returns false; otherwise it returns
unknown. The code does not import RuleIR or the reference evaluator. This makes
it useful for detecting a defect in how the expression was encoded or evaluated.

For this particular body, each input occurs once and with a single polarity.
The completion calculation agrees with the current three-valued operations.
That agreement should not be generalized to every expression. For example, a
formula containing a variable and its negation can be true under every Boolean
completion while a compositional three-valued evaluator returns unknown. The
next chapter derives that difference. The oracle's representation is independent
of the candidate AST, but its semantic scope remains specific.

Conflict, stale evidence and later-recorded evidence are not extra members of
the three-value enumeration. Their normalization and precedence are exercised
by named cases and other application checks. The count 729 must not be used to
imply exhaustive testing of those mechanisms. Nor does it exhaust numerical
inputs, dates, arbitrary rule graphs or the Java runtime's entire implementation.
A result summary should always name the finite object that was exhausted.

The comparison checked full semantic responses and relevant traces, while
respecting each backend's own identity. It would be wrong to require the Python
and Java engine identifiers to be identical: those identifiers describe
different implementations. It would also be wrong to ignore all metadata in a
comparison, because source, evidence and rule identities are part of an
explainable decision. Chapter~\ref{ch:verification} specifies the distinction
between semantic agreement and implementation-specific provenance.

\section{Mutations ask whether the tests can notice an error}
\label{sec:circular-mutations}
Successful cases can be uninformative if they would also pass a wrong program.
Mutation testing addresses this by deliberately changing the rule and asking
whether the checks detect the change. A mutation is not a plausible legal
interpretation merely because it can be compiled. It is an instrument for
testing the sensitivity of the evidence.

The first retained mutation removed the fee-discount exception. Cases g03, g07,
g16 and g19 detected the change. Some of these detections concern the logical
exception, while g19 also exposes the effect of changing the static dependency
closure. Removing an expression can remove a referenced fact from conflict
analysis. A mutation report should therefore explain why a case detects a
change rather than treating every detection as the same kind of evidence.

The second mutation omitted the product-type route. Cases g02, g14 and g15
detected it. Their relationship to the source is direct: they isolate the
route that the mutation removed, including a case where that route is
unresolved. A baseline case with a named fund could not provide this evidence.

The third mutation assumed that a gift existed without requiring its evidence.
Cases g06, g08, g18, g20, g21 and g22 detected the resulting change. The set
includes known absence, missing classification, conflict, expiry and later
knowledge. This shows why evidence-handling cases are necessary even for a
small logical condition. A wrong classification default can alter outcomes
across several superficially different scenarios.

The fourth mutation ignored the selected actor/product scope. Cases g10 through
g13 detected it. The expected difference includes both out-of-scope and unknown
results. If the test harness compared only a final allow/deny bit after
flattening all other statuses, it could conceal these errors.

All four mutations were compiled into Java before comparison. The result thus
concerns executable altered candidates, not a textual inspection of a proposed
patch. Still, detecting four mutations does not establish detection of every
possible implementation or interpretation defect. Mutation operators are chosen
by people or software and can themselves omit an important failure mode. The
ensemble evaluation must maintain a documented fault model and extend it when
new failures are discovered.

\section{A new discrepancy: what counts as a discount?}
\label{sec:circular-cashback}
The existing tests assume that a reviewer can supply the discount classification.
To investigate the missing interpretation step, consider a synthetic offer in
which a customer pays a fee and later receives cash back. Two candidate
definitions might be proposed. One restricts the exception to a reduction of
the fee initially charged. Another includes a documented refund of a fee that
was actually paid. The source packet alone has not established that both
readings are legally tenable. Their purpose here is to expose a question whose
answer affects implementation.

Fix the actor, activity and promotional link, and assume for the comparison
that the component otherwise meets the gift condition. Let the fee paid be
HK\$100 and the later rebate also HK\$100. Under the first proposed definition,
the later transfer is outside the defined exception. Under the second, it may
fall within the exception if the required refund relationship is established.
Those conditions must be derived from authority or review; the example does
not introduce them as new SFC requirements.

The next action should be a targeted investigation of the applicable definition,
Code and FAQ, not a request for the same model to vote again on the word
``discount.'' A source that explicitly includes the relevant refund treatment
would support refinement of the second candidate. A source that excludes it
would defeat that candidate in the applicable scope. A source discussing
another product, actor or period may be informative without settling the
question. The retrieval result needs an applicability assessment before it can
resolve the discrepancy.

Now alter the cash amount to HK\$150. If a candidate definition treats only a
refund of an actual HK\$100 fee as falling within the exception, the extra
HK\$50 raises a further classification question. It does not follow that the
system can automatically split the transfer into a lawful and unlawful part.
The legal unit of assessment and treatment of mixed benefits must first be
settled. An arithmetic decomposition is easy; its legal significance is the
question under investigation.

A second variation changes the evidence. The marketing material describes the
transfer as a fee refund, but no record establishes that a fee was paid. The
candidate's definition may be clear while its decisive factual premise remains
unknown. This discrepancy belongs to fact classification or evidence
availability, not necessarily to interpretation. The resolution controller
should request the relevant evidence or return unresolved, rather than
changing the definition to make the missing record unnecessary.

A third variation combines a refund with a separate voucher. A whole-offer
exception would remove the need to assess the voucher independently. The
proposed component-level interpretation would retain it. This is a different
branch in the hypothesis space from the timing of a refund. Preserving separate
issue identifiers prevents a source that resolves the timing question from
being incorrectly treated as resolving the mixed-offer question as well.

These variations show how an ensemble can make review more precise. It should
produce a small set of materially different definitions, the source reasons
for each, and the smallest scenarios on which they diverge. The reviewer then
decides an identifiable question. A long narrative saying that an offer appears
reasonable would provide less actionable evidence.

\section{A complete dry run of the proposed workflow}
\label{sec:circular-dryrun}
The following dry run specifies proposed behavior. It is not a report that the
new ensemble has already run. Start by creating an assessment whose context
contains the retained circular revision, declared distributor/fund scope,
selected effective interval and an inventory of dependencies. The historical
engineering example uses a synthetic September 2026 interval; the workflow must
not infer a regulatory commencement date from that setting.

The source-review path inventories the circular and flags the incorporated Code
and FAQ as unresolved. The interpretation path proposes the component-level
gift rule. A second interpretation path omits the product-type route. A
definition-focused path asks whether cashback is a discount and whether a
combined offer changes the assessed unit. At this stage the system has a
candidate set and several issues, not a winner.

The comparator identifies a decision difference between the complete promotional
condition and the candidate omitting the product-type route. It constructs a
case with a gift linked only to a fund category and no exception. The
source-grounding check points to the relevant paragraph. The candidate author
can propose a repair adding the missing route. The system stores the repaired
candidate as a new revision, reruns the discrepancy case and the retained
regression cases, and preserves the original candidate and its rejection reason.

The cashback definition issue follows a different route. Retrieval attempts to
obtain the incorporated sources with the right identity and version. If the
sources are unavailable, another language-model explanation cannot close the
issue. The controller records the failed acquisition and the remaining
dependency. It may retry an authorized retrieval route within its budget, or
stop automatic work because no permitted action can supply the missing
authority. The issue remains unresolved either way.

The mixed-offer issue can produce a useful witness even before legal resolution.
The comparison holds the voucher facts fixed and toggles the presence of a
separate fee waiver. If a candidate's voucher decision changes solely because
of the other component's exception, the report explains the dependence. A
reviewer must then confirm the unit to which the exception applies. The witness
does not grant the controller authority to make that legal decision itself.

The accepted formal candidate can be compiled and checked in parallel with this
research. Candidate compilation is valuable because it exposes implementation
defects early. It does not change the issue status. The report may legitimately
say that the candidate passed all specified execution checks and remains
unavailable for release because source interpretation is unresolved. Keeping
those two observations together is more informative than suppressing either.

When automatic work reaches its finite limit, the controller freezes a report.
It identifies the source revision, surviving definitions, rejected encodings,
discriminating scenarios, completed actions, failed actions, unresolved
dependencies and the exact effect of each unresolved issue on the control.
There is no final vote to select a legally preferred answer. The state is ready
for targeted human review or blocked on a named dependency, according to the
reason for stopping.

A designated reviewer can later resolve an issue using newly obtained authority.
That creates a new context revision. Dependent candidate scores and comparisons
are invalidated where appropriate, and the formal rule and tests are updated.
The earlier report remains a faithful account of what was unresolved before the
new material arrived. This is essential for reproducibility: a current reviewer
must be able to see why the earlier system abstained.

\section{From the JSON oracle to independent adjudication}
\label{sec:circular-independent}
The filename \texttt{oracle.json} is an engineering term. It designates a source
of expected results for tests. It does not establish institutional authority or
legal independence. The file states that the same agent authored it and the
candidate. A future user interface should display that authorship rather than
letting the label ``oracle'' imply a gold standard.

Independent adjudication should begin from the retained source packet and
declared task, before exposing reviewers to the proposed formal expression.
One reviewer identifies requirements, definitions and difficult cases. Another
can independently develop or challenge the interpretation. Disagreements are
recorded and, where possible, resolved with source reasons. A third role may
adjudicate material disagreement. Independence here concerns information and
responsibility as well as the identity of a person or model.

After that source-based stage, reviewers can inspect the candidate rules and
their discriminating cases. This order reduces anchoring on the generated
answer. It does not make reviewers infallible. The reference packet must allow
an unresolved category, later correction and competing defensible readings.
If a reviewer cannot determine the treatment of a cash rebate without the
referenced FAQ, the correct reference outcome is not a guessed Boolean chosen
to make benchmark scoring convenient.

The original 22 cases should remain frozen as historical engineering evidence.
New adjudicated cases belong in a separately versioned packet with its own
authors, source context and review decisions. A changed expected outcome should
explain whether the earlier interpretation was rejected, the scope changed,
new evidence became available or the evaluation semantics were corrected.
Those causes imply different repairs and different consequences for existing
releases.

The chapter has now followed a source to a rule, tested the rule, exposed
unresolved definitions and specified how a future ensemble should handle them.
The remaining task is to state precisely what each formal check can prove.
Without that precision, a solver's success could once again be mistaken for
the missing judgment about meaning.
